\begin{ouverture}[college1]{1}{Masse, volume et masse volumique}{Partie 1 · Organisation et transformations de la matière}{Comment reconnaître un matériau sans le découper ?}
\objectif{Mesurer une masse et un volume avec les unités adaptées.}
\objectif{Construire puis exploiter la relation entre masse et volume.}
\objectif{Identifier les limites d’une comparaison expérimentale.}
\prerequis{Lire une graduation et utiliser la proportionnalité.}
\end{ouverture}
\pcPage{Le cours}{Masse, volume et masse volumique}
\phantomsection\label{ess-1}
\essSection{1}{Mesurer deux grandeurs différentes}
La masse $m$ mesure une quantité de matière au sens usuel du collège ; on la mesure avec une balance. Le volume $V$ est l’espace occupé. Pour un liquide, placer l’œil à hauteur du bas du ménisque d’une éprouvette graduée. Tarer la balance avec le récipient vide permet de mesurer directement la masse du liquide. On note toujours valeur et unité.\par
\pcFigure{figures/c01-mesure.pdf}{Un objet immergé déplace un volume d’eau égal à son volume, s’il est totalement immergé et ne se dissout pas.}
\essSection{2}{Passer d’une unité à une autre}
$1\unite{L}=1\unite{dm^3}$, $1\unite{mL}=1\unite{cm^3}$ et $1\unite{m^3}=1\,000\unite{L}$. Les conversions cubiques demandent de la vigilance : $1\unite{m^3}=1\,000\,000\unite{cm^3}$. Pour un solide compact insoluble, la différence entre les volumes final et initial de l’eau donne le volume déplacé. Une bulle collée au solide fausse ce résultat.\par
\clearpage
\essSection{3}{Une propriété du matériau}
Pour une substance homogène, à température et pression fixées, masse et volume sont proportionnels. Le quotient $\rho=m/V$, appelé masse volumique, caractérise la substance dans ces conditions. On utilise aussi $m=\rho V$ et $V=m/\rho$. Avec $m$ en g et $V$ en mL, $\rho$ s’exprime en g/mL. $1\unite{g/mL}=1\,000\unite{kg/m^3}$.\par
\pcFigure{figures/c01-rho.pdf}{Les pentes des droites diffèrent : un même volume de deux substances peut avoir des masses différentes.}
\essSection{4}{Comparer avec prudence}
À environ $20\,^{\circ}\mathrm C$, on prend ici $\rho_{\text{eau}}\approx1{,}00\unite{g/mL}$, $\rho_{\text{aluminium}}\approx2{,}70\unite{g/cm^3}$ et $\rho_{\text{fer}}\approx7{,}9\unite{g/cm^3}$. Une valeur proche aide à identifier un matériau, sans suffire à prouver sa composition. La température, une cavité ou un mélange peuvent modifier la mesure. Chauffer un liquide le dilate généralement : sa masse volumique diminue si sa masse reste constante.\par
\begin{pcEncadre}{method}{Un exemple expliqué}Un solide a une masse de $54{,}0\unite g$. L’éprouvette passe de $30{,}0$ à $50{,}0\unite{mL}$. Son volume vaut $50{,}0-30{,}0=20{,}0\unite{cm^3}$ ; donc $\rho=54{,}0/20{,}0=2{,}70\unite{g/cm^3}$. La valeur est compatible avec de l’aluminium compact. On vérifie l’immersion et l’absence de bulles avant de conclure.\end{pcEncadre}
\pcPage{Expérimenter et retenir}{Masse, volume et masse volumique}
\phantomsection\label{activite-1}
\begin{pcEncadre}{method}{Investigation en classe}\textbf{Question.} Un liquide inconnu est-il compatible avec de l’eau ? En classe, peser un récipient sec, puis y verser successivement trois volumes connus du même liquide, sans changer la température. Mesurer chaque masse par différence. Construire $m$ en fonction de $V$, avec des axes gradués et des unités. Comparer les quotients et leur dispersion. Proposer une explication si un point s’écarte. Ne pas goûter. \textbf{À rendre :} tableau, graphique, calcul et conclusion limitée à la précision des instruments.\end{pcEncadre}
\pcSousTitre{Ce que je dois savoir expliquer}\begin{questions}
\item La masse se mesure avec une balance ; le volume avec une méthode adaptée.
\item $V_{\text{solide}}=V_f-V_i$ pour une immersion convenable.
\item $\rho=m/V$ avec des unités cohérentes.
\item À conditions fixées, doubler le volume d’une substance double sa masse.
\item Une identification exige une comparaison et une discussion des limites.
\end{questions}
\begin{pcEncadre}{attention}{Une erreur à éviter}« Le plus lourd est le plus dense. » Une grande bouteille d’huile peut être plus lourde qu’un petit verre d’eau tout en ayant une masse volumique plus faible.\end{pcEncadre}
\pcRepere{cor-activite-1}{Pistes pour l’investigation}\pcGrille
\clearpage\markright{Je démarre}
\rubrique{Je démarre}
\par\noindent\begin{minipage}{\linewidth}
\begin{exercice}[1]{Choisir l’instrument}
\begin{questions}
\item Quel appareil mesure la masse d’un caillou ?
\item Quelle grandeur exprime une indication en mL ?
\end{questions}
\end{exercice}
\end{minipage}
\par\noindent\begin{minipage}{\linewidth}
\begin{exercice}[1]{Convertir sans confondre}
\begin{questions}
\item Convertir $0{,}250\unite L$ en mL.
\item Convertir $35\unite{cm^3}$ en mL.
\end{questions}
\end{exercice}
\end{minipage}
\par\noindent\begin{minipage}{\linewidth}
\begin{exercice}[1]{Lire une immersion}
\pcFigureEx{figures/ex-immersion.pdf}{Lecture de deux volumes ; solide entièrement immergé.}
\begin{questions}
\item L’eau passe de $42$ à $57\unite{mL}$. Quel est le volume du solide ?
\item Nommer une condition nécessaire.
\end{questions}
\end{exercice}
\end{minipage}
\par\noindent\begin{minipage}{\linewidth}
\begin{exercice}[1]{Une tare utile}
\begin{questions}
\item Le récipient vide pèse $82\unite g$ ; rempli, $132\unite g$. Calculer la masse du liquide.
\item Que ferait la touche tare avant le remplissage ?
\end{questions}
\end{exercice}
\end{minipage}
\par\noindent\begin{minipage}{\linewidth}
\begin{exercice}[1]{Calcul direct}
\begin{questions}
\item $40\unite{mL}$ d’un liquide ont une masse de $32\unite g$. Calculer sa masse volumique.
\item Est-elle supérieure à celle de l’eau ?
\end{questions}
\end{exercice}
\end{minipage}
\clearpage\markright{Je m’entraîne}
\rubrique{Je m’entraîne}
\par\noindent\begin{minipage}{\linewidth}
\begin{exercice}[2]{Une masse à prévoir}
\begin{questions}
\item Un matériau a $\rho=2{,}7\unite{g/cm^3}$. Quelle masse pour $30\unite{cm^3}$ ?
\item Quelle masse pour un volume double ?
\end{questions}
\end{exercice}
\end{minipage}
\par\noindent\begin{minipage}{\linewidth}
\begin{exercice}[2]{Trouver le volume}
\begin{questions}
\item Un objet compact de $79\unite g$ a $\rho=7{,}9\unite{g/cm^3}$. Trouver $V$.
\item Exprimer ce résultat en mL.
\end{questions}
\end{exercice}
\end{minipage}
\par\noindent\begin{minipage}{\linewidth}
\begin{exercice}[2]{Reconnaître une proportionnalité}
\pcFigureEx{figures/ex-rho-points.pdf}{}
\begin{questions}
\item Pour les volumes $10$, $20$ et $30\unite{mL}$, on mesure $9$, $18$ et $27\unite g$. Le modèle proportionnel convient-il ?
\item Quelle valeur de $\rho$ en déduire ?
\end{questions}
\end{exercice}
\end{minipage}
\par\noindent\begin{minipage}{\linewidth}
\begin{exercice}[2]{Comparer à volume égal}
\begin{questions}
\item Deux blocs de même volume $12\unite{cm^3}$ sont en aluminium et en fer. Calculer leurs masses avec les données du cours.
\item Lequel a la plus grande masse volumique ?
\end{questions}
\end{exercice}
\end{minipage}
\par\noindent\begin{minipage}{\linewidth}
\begin{exercice}[2]{Changement d’unité}
\begin{questions}
\item Exprimer $0{,}85\unite{g/mL}$ en kg/m$^3$.
\item Quelle masse ont $2\unite L$ de ce liquide ?
\end{questions}
\end{exercice}
\end{minipage}
\clearpage\markright{Je raisonne}
\rubrique{Je raisonne}
\par\noindent\begin{minipage}{\linewidth}
\begin{exercice}[3]{Une bulle gênante}
\begin{questions}
\item Une bulle adhère au solide immergé. Le volume mesuré est-il trop grand ou trop petit ?
\item Quel effet cela a-t-il sur $m/V$ ?
\end{questions}
\end{exercice}
\end{minipage}
\par\noindent\begin{minipage}{\linewidth}
\begin{exercice}[3]{Un objet creux}
\begin{questions}
\item Une pièce de $27\unite g$ occupe extérieurement $20\unite{cm^3}$. Calculer sa masse volumique moyenne.
\item Cette valeur suffit-elle à exclure une coque en aluminium ?
\end{questions}
\end{exercice}
\end{minipage}
\par\noindent\begin{minipage}{\linewidth}
\begin{exercice}[3]{Diagnostiquer un graphique}
\pcFigureEx{figures/ex-tare-graphe.pdf}{Exemple de données illustrant une pesée sans tare.}
\begin{questions}
\item Un élève pèse le récipient avec chaque quantité de liquide sans soustraire sa masse. La droite passe-t-elle par l’origine ?
\item Comment corriger les données ?
\end{questions}
\end{exercice}
\end{minipage}
\par\noindent\begin{minipage}{\linewidth}
\begin{exercice}[3]{Comparer deux protocoles}
\begin{questions}
\item Une éprouvette est graduée tous les $2\unite{mL}$. On peut immerger une bille ou dix billes identiques à la fois. Quel protocole réduit l’effet relatif de la lecture ?
\item On trouve $40\unite{mL}$ pour dix billes. Calculer le volume moyen.
\end{questions}
\end{exercice}
\end{minipage}
\par\noindent\begin{minipage}{\linewidth}
\begin{exercice}[3]{Un liquide qui se dilate}
\begin{questions}
\item La masse reste $200\unite g$ ; le volume passe de $200$ à $204\unite{mL}$. Calculer les deux masses volumiques.
\item Expliquer le sens de la variation.
\end{questions}
\end{exercice}
\end{minipage}
\pcPage{Défi du chapitre}{Masse, volume et masse volumique}
\phantomsection\label{defi-ex-1}
\pcDefiTitre{1}{La pièce mystérieuse}
\begin{questions}
\item Une pièce pèse $43{,}2\unite g$ et déplace $16{,}0\unite{mL}$. Calculer $\rho$.
\item Une référence donne $2{,}70\unite{g/cm^3}$ pour l’aluminium. Formuler une conclusion.
\item Proposer deux contrôles avant d’identifier le matériau.
\end{questions}
\pcGrille
